%% Cover letter (draft, for the authors to edit and sign): Paper A to Annals of Global Analysis and
%% Geometry. Build: pdflatex cover-letter-A-AGAG.tex
\documentclass[11pt]{article}
\usepackage[a4paper,top=16mm,bottom=16mm,left=20mm,right=20mm]{geometry}
\usepackage{amsmath,amssymb}
\usepackage[hidelinks]{hyperref}
\setlength{\parindent}{0pt}
\setlength{\parskip}{4pt}
\sloppy
\pagestyle{empty}

\begin{document}

\begin{flushright}
Palaash Gang\\
Indus International School, Pune, India\\
\texttt{palaash.gang@indusschoolpune.com}\\[6pt]
\today
\end{flushright}

The Editors-in-Chief\\
\emph{Annals of Global Analysis and Geometry}

Dear Editors,

On behalf of all six authors (Palaash Gang, Akshaj Agadi, Arjun Veluri, Jerry Wang, Jeremy Barreto and Aarin Chouthaiwale), I submit the manuscript \emph{How much of a hyperbolic orbifold does heat hear?} for consideration in the \emph{Annals of Global Analysis and Geometry}. It is posted as arXiv:[PLACEHOLDER: arXiv number of Paper A], and it has an electronic supplement (Online Resource~1).

\textbf{What the paper does.} The heat trace of a closed hyperbolic $2$-orbifold has a short-time expansion whose coefficients, the heat invariants, are spectral invariants. The paper asks how many heat invariants are needed to determine the signature, that is, the genus and the orders of the cone points. If every cone order is at most $M$, the first $M+1$ heat invariants suffice whatever the area, and $M+1$ is attained, so the number can grow only through large cone orders. Without that bound, the number needed among orbifolds of area at most $A$ grows at least like $\sqrt A$ and at most linearly, and it grows at least like a power $A^\alpha$ exactly when the Prouhet--Tarry--Escott function satisfies $N(k)=O(k^{1/\alpha})$; linear growth is equivalent to the open problem $N(k)=O(k)$. Among triangle orbifolds three heat invariants always suffice, and two first fail on $\mathcal O(2,8,8)$ and $\mathcal O(3,3,12)$, a pair isolated by an elliptic curve of rank~$0$. The paper also shows that this pair carries metrics, flat near the cone points, whose heat expansions agree to all orders, so constant curvature is essential.

\textbf{Why AGAG.} The paper continues a line of work on heat invariants and spectral data of orbifolds that the journal has published: Stanhope's spectral bounds on orbifold isotropy (\emph{Ann.\ Global Anal.\ Geom.}\ 27, 2005), Abreu, Dryden, Freitas and Godinho on hearing the weights of weighted projective planes (33, 2008), Rossetti, Schueth and Weilandt on isospectral orbifolds with different maximal isotropy orders (34, 2008), and Schueth's heat coefficients of surfaces with curved conical singularities (69, 2026). The paper cites each of these.

\textbf{Companion manuscripts.} Two companion manuscripts by the same authors are being submitted at the same time.
\begin{itemize}\setlength{\itemsep}{2pt}
\item \emph{Finitely many eigenvalues determine the signature of a hyperbolic orbifold}, submitted to the \emph{Annals of Global Analysis and Geometry} (arXiv:[PLACEHOLDER: arXiv number of Paper B]). It shows that finitely many approximate eigenvalues determine the signature, with explicit bounds. It uses results of Sections~2 and~3 of the present paper, with this paper's proofs.
\item \emph{Triples with equal sum and equal reciprocal sum}, submitted to \emph{Integers} (arXiv:[PLACEHOLDER: arXiv number of the arithmetic note]). It studies the arithmetic of the triples of cone orders that two heat invariants do not distinguish.
\end{itemize}
No proof in the present paper depends on either companion. The three manuscripts share one code and data deposit, and the present paper shares its computed spectra with the first companion.

\textbf{Code and data.} The code and data are publicly archived at Zenodo, with all six authors as creators: [PLACEHOLDER: Zenodo DOI, to be minted at submission]. One command reruns every check. The heat invariants, certificates and search results are exact unless stated otherwise. The supplement lists the statements that rest on computer search alone, and none of them is used in a proof.

The manuscript is not under consideration elsewhere, and all authors have approved its submission. We would be grateful if you would consider it for publication.

Yours sincerely,

\vspace{7mm}
Palaash Gang (corresponding author), on behalf of all authors

\end{document}
